% Extension to insert after hibridacion.tex and the thermal module.
% The complete electronic equation and towers remain in their own sections.
\section{Electronic History and Mass Conservation under Transport}
\label{md:r27:electron}

Equation \eqref{md:compos:electron-completo} evaluates a persistent central
section and preserves its action return. The trajectory, its memory and
its mass are distinct readings of the same APP--TRIT--TPK state.
Transport of a route may modify events and memory while preserving the mass
reader's value. This is decided using the complete character of
\eqref{md:hib:caracter}, together with the return exponent.

\subsection{Separation and Reconstruction of the Two Central Histories}
The expressed event registers of the two orientations are
\begin{equation}
 e_{++}=(2,2,5,-4,-3,-2)^2,\qquad
 e_{--}=(4,3,2,-2,-2,-5)^2,
 \label{md:r27:historias}
\end{equation}
where the exponent $2$ denotes concatenation of two blocks. They have zero
total and the same absolute histogram. For a register $e\in\mathbb Z^{12}$,
let
\begin{equation}
 \begin{gathered}
 p_i=e_i+e_{i+6},\quad \mathfrak a_i=e_i-e_{i+6},\\
 s_C=\sum_{i=0}^{5}p_i,\quad M_i=p_i-p_5\quad(0\le i<5).
 \end{gathered}
 \label{md:r27:registro-electron}
\end{equation}
\begin{proposition}[Integral reconstruction of the register]
The data $(s_C,M_0,\ldots,M_4,\mathfrak a_0,\ldots,\mathfrak a_5)$ determine the register through
\begin{equation}
 \begin{gathered}
 p_5=\frac{s_C-\sum_{i=0}^{4}M_i}{6},\quad p_i=p_5+M_i,\\
 e_i=\frac{p_i+\mathfrak a_i}{2},\quad e_{i+6}=\frac{p_i-\mathfrak a_i}{2}.
 \end{gathered}
 \label{md:r27:inversa-electron}
\end{equation}
Its integral image is characterized by divisibility of the first numerator
by six and the congruence $p_i\equiv \mathfrak a_i\pmod2$ for each $i$.
\end{proposition}
\begin{proof}
Summing $p_i=p_5+M_i$ over the first five indices and adding $p_5$ gives
$s_C=6p_5+\sum M_i$. Adding and subtracting the definitions of $p_i,\mathfrak a_i$
gives the remaining formulas. The stated conditions are necessary for integer
values and sufficient to reconstruct them; substitution recovers all the
original data.
\end{proof}
For \eqref{md:r27:historias}, $\mathfrak a=0$ and
$M_{++}=(8,8,14,-4,-2)$, whereas $M_{--}=(18,16,14,6,6)$.
Reconstruction distinguishes both histories even when particular scalar
readers share their value. The symbols $\mathfrak a_i$ in this decomposition are event
differences and are distinguished from the areal character $a_0$ by their domain.

\subsection{Spectral Form of the Return Reader}
For $\tau\in\mathbb R^{12}$, with cyclic indices, define
\begin{equation}
 C_s=\sum_{j=0}^{11}\tau_j\tau_{j+s},\quad
 T_k=\sum_{j=0}^{11}\tau_j e^{-2\pi i kj/12},\quad
 \theta_k=2\pi k/12.
 \label{md:r27:fourier}
\end{equation}
The character representation is applied after the register has been produced.
The return reader retains
$\Omega=-2C_1-4C_2-6C_3$ and $P_6=C_6/2$.
\begin{proposition}[Exact harmonic decomposition]
\begin{equation}
 \begin{split}
 C_s&=\frac1{12}\sum_{k=0}^{11}|T_k|^2\cos(s\theta_k),\\
 \Omega&=\frac1{12}\sum_{k=0}^{11}|T_k|^2
 [-2\cos\theta_k-4\cos2\theta_k-6\cos3\theta_k],\\
 P_6&=\frac1{24}\sum_{k=0}^{11}|T_k|^2(-1)^k.
 \end{split}
 \label{md:r27:espectro}
\end{equation}
\end{proposition}
\begin{proof}
Expanding $|T_k|^2$ gives
$\sum_{j,l}\tau_j\tau_l e^{-i\theta_k(j-l)}$.
The sum $\sum_{k=0}^{11}e^{i\theta_km}$ is twelve when $12$ divides $m$
and zero otherwise. Applying this identity to the shift $s$ recovers $C_s$.
Since $\tau$ is real, imaginary parts cancel in pairs $k,12-k$, leaving the
cosine. The other two formulas follow by linear combination and by
$\cos(6\theta_k)=(-1)^k$.
\end{proof}
For $\tau_e=(+++---)^2$, the only nonzero powers are
$|T_2|^2=64$, $|T_6|^2=16$, $|T_{10}|^2=64$. Their weights in $\Omega$
are $7,4,7$, so that $\Omega=80$ and $P_6=6$.
This recovers the harmonic part of
$\kappa_e=80-54\alpha+6\Delta_4$ in the complete electronic equation.
The alternating register $(+-)^6$, with the same histogram, has
$|T_6|^2=144$ and $\Omega=48$: order is an effective dependency of the reader.
Cyclic translations preserve $|T_k|^2$ while changing the phase of $T_k$;
this spectral reading preserves correlations, rather than the entire history.

\subsection{Criterion for Transported Mass and Its Associated Readings}
Let $\mathcal H_{\mathcal F}$ be the fibre of a finite route domain, with
basis $\{|\Gamma\rangle\}$. Fix one return reader, $r=D$ or $r=\Omega$,
and a common chart with $R_{\rm act}>0$.
The mass law of \eqref{md:hib:factor-ruta} determines
\begin{equation}
 \begin{gathered}
 \frac{m_\Gamma^\beta}{m_e^\beta}
 =\chi_{\rm ref}(\sigma_\Gamma-\sigma_e,\eta_\Gamma-\eta_e)
 R_{\rm act}^{\kappa_\Gamma-\kappa_e},\\
 \Lambda_\Gamma=\ln(m_\Gamma^\beta/m_e^\beta).
 \end{gathered}
 \label{md:r27:masa-relativa}
\end{equation}
The logarithm acts on a dimensionless ratio. On a nonscalar fibre, the
operator and its declared diagonal realization are retained; the finite domain
may be an orbit or multisection, without choosing an arbitrary individual cell.

\begin{theorem}[Conservation of the mass reader on a stable domain]
Let $F$ be a TPK transport acting bijectively on the fixed domain, and let
$V|\Gamma\rangle=|F\Gamma\rangle$. For
$M|\Gamma\rangle=m_e^\beta e^{\Lambda_\Gamma}|\Gamma\rangle$,
\begin{equation}
 [M,V]|\Gamma\rangle=m_e^\beta
 (e^{\Lambda_{F\Gamma}}-e^{\Lambda_\Gamma})|F\Gamma\rangle.
 \label{md:r27:conmutador-masa}
\end{equation}
Thus $[M,V]=0$ is equivalent to $\Lambda_{F\Gamma}=\Lambda_\Gamma$ for
every route. With chart parameters fixed, the condition is
\begin{equation}
 \begin{aligned}
 &\delta n_A x+\frac{\delta s_C}{6}y+
 \delta k_{\Delta_4}\Delta_4+\delta b_\zeta\zeta_{270}\\
 &\qquad+\sum_h\delta N_h(q_-^h-q_+^h)
 +\delta\kappa\ln R_{\rm act}=0,
 \end{aligned}
 \label{md:r27:criterio-masa}
\end{equation}
where $\delta$ compares the transported route with the initial one.
\end{theorem}
\begin{proof}
Applying $MV$ and $VM$ to every basis vector gives
\eqref{md:r27:conmutador-masa}. Positivity of $m_e^\beta$ and injectivity
of the real exponential give the equivalence. Taking the logarithm of
\eqref{md:r27:masa-relativa} and subtracting the two route characters gives
\eqref{md:r27:criterio-masa}. For infinite towers, require for each route
$\sum_h|N_h(q_-^h-q_+^h)|<\infty$; absolute convergence legitimizes
the subtraction. Finiteness of the route domain preserves the bounded
operators used in this proof.
\end{proof}

Transport is first produced by TPK; the identity subsequently determines
what it preserves. Individual events and memory may change and compensate
in \eqref{md:r27:criterio-masa}. With the action section, clock and circular
realization of \eqref{md:hib:reciprocidad-radial} fixed, commutation with
$M$ also preserves the proper-frequency, gravitational-radius, reduced-Compton-
wavelength and thermal-weight readings that are functions of the same positive
operator. Inversion of mass is restricted to the complement of the null sector.
In the marked realization of \eqref{md:r27:area-gravedad}, the fourth reading is
\begin{equation}
 \Xi_\Gamma=\frac{\Omega_\Gamma}{\omega_0}
 =\frac{r_{g,\Gamma}}{L_\alpha}
 =\frac{L_\alpha}{\bar\lambda_\Gamma}
 =\frac{E_\Gamma}{b_*\Theta_{{\rm clk},P}\ln3}.
 \label{md:r27:lecturas-conservadas}
\end{equation}
The proof substitutes the identities already established in a common chart.
This criterion describes conservation under a declared transport; the physical
selectors of species, representation and interaction retain their own
conditions. The trajectory and its harmonic reading add information to the
scalar mass without replacing those conditions.
